% !TEX root = ../mmskills_neurips.tex

\section{Detailed Related Work}
\label{app:related-work}

This section provides the expanded related-work discussion summarized in Section~\ref{sec:related-work}.

\paragraph{Skills for agents.}
Skill reuse has a long history in temporal abstraction for reinforcement learning and motor primitives for robotics \citep{sutton1999options,ijspeert2013dmp}. Recent LLM agents have made skills a practical interface for storing and composing procedural knowledge in language-conditioned environments. Early systems connected language models to action by grounding language in affordances \citep{ichter2022saycan}, emitting executable programs \citep{liang2023codepolicies}, or interleaving reasoning and acting \citep{yao2023react}; adjacent code- and tool-agent work studies robust tool-call data loops, search-based code refinement, and adversarial test-case generation \citep{zhang2025looptoolclosingdatatrainingloop,li2025rethinkmctsrefiningerroneousthoughts,li2026atgen}. Reflection mechanisms then made agent behavior more persistent across attempts \citep{shinn2023reflexion}. In open-ended environments, systems such as DEPS, Voyager, and JARVIS-1 showed that large models can use language, stored experience, and self-generated programs to acquire or reuse behaviors over extended task horizons \citep{wang2024deps,wang2023voyager,wang2023jarvis}. These works motivate our focus on procedural reuse, but their reusable knowledge is primarily textual, symbolic, or programmatic.

More recent work treats skills as an explicit substrate for agent improvement. SkillWeaver distills web exploration into reusable API-like skills \citep{zheng2025skillweaver}; CUA-Skill builds a parameterized skill base with execution and composition graphs for computer-using agents \citep{chen2026cuaskill}; SkillX automatically constructs hierarchical skill knowledge bases from agent experience \citep{wang2026skillx}; EvoSkill studies automated skill discovery through failure analysis in multi-agent settings \citep{alzubi2026evoskill}, where decentralized coordination and scalable improvement are also central concerns \citep{yang2025agentnetdecentralizedevolutionarycoordination,shao2026monoscalescalingmultiagentmonotonic}; SkillClaw evolves shared skills from multi-user trajectories \citep{ma2026skillclawletskillsevolve}; and SkillRL co-evolves a hierarchical skill library with reinforcement learning \citep{xia2026skillrlevolvingagentsrecursive}. A recent survey frames agent skills as portable packages of instructions, code, and resources loaded through progressive disclosure \citep{xu2026agentskills}. A complementary perspective treats accumulated agent experience as long-term memory: Generative Agents maintain a memory stream that supports recall, reflection, and planning \citep{park2023generativeagentsinteractivesimulacra}, while MemGPT introduces an OS-style memory hierarchy that pages information in and out of the model's working context \citep{packer2024memgptllmsoperatingsystems}. MMSkills follows this broader move toward modular procedural knowledge, but changes the unit being stored: instead of treating skills mainly as text, code, APIs, or execution graphs, we define a skill package whose central evidence is a set of visually grounded runtime states. Branch loading also takes inspiration from memory-paging ideas, by inspecting selected multimodal evidence in a temporary branch rather than flooding the main context.

This emerging ecosystem has also motivated dedicated evaluation of skill utility. SkillsBench measures how skills affect agent performance across diverse tasks \citep{li2026skillsbenchbenchmarkingagentskills}, SkillTester evaluates utility and security risks of agent skills \citep{wang2026skilltesterbenchmarkingutilitysecurity}, and recent work studies skill usage under more realistic retrieval and adaptation settings \citep{liu2026agenticskillsworkwild}. These benchmarks show that skills are not automatically beneficial; their value depends on relevance, compactness, selection, and safe use, especially as self-evolving agents may introduce emergent risks \citep{shao2026agentmisevolveemergentrisks}. Our work addresses a complementary question for visual agents: what evidence should a skill expose, and how should that evidence be loaded, when correct use depends on the current visual state?

The closest line to our work is multimodal and GUI-specific skill augmentation. Mirage-1 introduces hierarchical multimodal skills for GUI agents and uses them with search to support long-horizon control \citep{xie2025mirage}; XSkill continually extracts experiences and skills for multimodal agents from visually grounded rollouts \citep{jiang2026xskillcontinuallearningexperience}; MuSEAgent studies stateful experiences for multimodal reasoning agents \citep{wang2026museagentmultimodalreasoningagent}; and CUA-Skill builds computer-use skills as parameterized procedures and execution graphs \citep{chen2026cuaskill}. MMSkills differs in emphasis: we define the skill artifact around reusable visual state evidence, not only around executable procedure structure or memory accumulation. Each skill is organized around when-to-use conditions, visible cues, verification cues, and multi-view state evidence, and the runtime first selects the relevant evidence before exposing it to the main agent. This makes the contribution a representation and loading mechanism for multimodal procedural cues, rather than another text skill library or GUI action graph.

\paragraph{Visual agents.}
Visual agents have rapidly advanced from web navigation to general computer use. Benchmarks such as Mind2Web and WebArena established realistic web-agent evaluation beyond synthetic interfaces \citep{deng2023mind2web,zhou2024webarena}; VisualWebArena showed that many web tasks require visual grounding rather than text-only reasoning \citep{koh2024visualwebarena}; and WebVoyager demonstrated end-to-end web interaction with large multimodal models on real websites \citep{he2024webvoyager}. The same trend appears in mobile, desktop, and embodied settings: Android in the Wild and AndroidWorld study device control from visual UI observations \citep{rawles2023androidwild,rawles2025androidworld}, OSWorld and macOSWorld evaluate agents in real operating-system environments \citep{xie2024osworld,yang2025macosworld}, RiOSWorld evaluates risks in multimodal computer-use agents \citep{yang2025riosworldbenchmarkingriskmultimodal}, and VisualAgentBench includes VAB-Minecraft and VAB-OmniGibson for open-world and household embodied interaction \citep{liu2024visualagentbench}.

Model and framework work has likewise moved toward visually grounded action, reflecting the shared multimodal objective of aligning visual and textual representations \citep{liu2024alignrecaligningtrainingmultimodalrecommendations,zhang2024dreamdualrepresentationlearningmodel}. SeeClick trains GUI grounding for screenshot-only agents \citep{cheng2024seeclick}; CogAgent introduces a visual language model dedicated to GUI understanding and operation \citep{hong2024cogagentvisuallanguagemodel}; OS-ATLAS learns a foundation action model for GUI control \citep{wu2024osatlas}; UI-TARS develops native GUI agents that perceive screenshots and emit keyboard/mouse actions \citep{qin2025uitars}; SeeAct builds web agents around general-purpose vision-language models \citep{zheng2024gpt4visiongeneralistwebagent}; AppAgent learns smartphone skills from on-device demonstrations \citep{zhang2023appagentmultimodalagentssmartphone}; OmniParser provides a pure-vision parser that turns screenshots into structured GUI elements \citep{lu2024omniparserpurevisionbased}; and Agent S provides a general computer-use framework built around GUI interaction \citep{agashe2024agents}. These systems improve the agent's perceptual and action interface. MMSkills instead targets the external knowledge layer used by such agents. A stronger GUI action model may click more accurately, but it still benefits from knowing which procedural state matters, which visual cue confirms progress, and which state indicates that a skill should not be applied. MMSkills represents that knowledge as a compact, reusable multimodal skill package.

\paragraph{GUI grounding benchmarks.}
Alongside task-completion benchmarks, a separate line of work measures how reliably GUI agents can localize UI elements from natural-language instructions. ScreenSpot-Pro extends earlier ScreenSpot evaluations to high-resolution, professional desktop environments, where target elements often occupy less than $0.1\%$ of the screen and the strongest grounding models still fall well below human performance \citep{li2025screenspotproguigroundingprofessional}. \citet{gou2025navigatingdigitalworldhumans} push toward universal visual grounding that lets agents identify GUI elements purely from screenshots, in the spirit of how humans navigate digital interfaces. MMBench-GUI organizes evaluation hierarchically, from content understanding and element grounding to task automation and multi-agent collaboration \citep{wang2025mmbenchguihierarchicalmultiplatformevaluation}, and DeskVision contributes a large-scale desktop dataset and evaluation suite that broadens grounding research across operating systems \citep{xu2025deskvisionlargescaledesktop}. These benchmarks isolate the perceptual layer of visual agents. MMSkills is complementary: rather than improving where to click, it provides procedural and visual evidence about which state matters at each step, and lets the underlying grounding capability translate that evidence into precise actions.

\paragraph{Long-context reliability.}
Recent studies have shown that simply enlarging the context window does not guarantee that all evidence is used effectively. \citet{liu2023lostmiddlelanguagemodels} report that language models often fail to retrieve information placed in the middle of long contexts, and benchmarks such as LongBench reveal substantial degradation as the input grows in length and modality \citep{bai2024longbenchbilingualmultitaskbenchmark}. These observations motivate our branch-loaded design: rather than directly inserting state cards, multi-view keyframes, and transition examples into the main agent context, the runtime first inspects selected evidence in a temporary branch and returns a compact structured guidance tuple. This isolates expensive multimodal evidence reading from action generation, and avoids the long-context failure modes that arise when reference views and live observations compete for the same context window.
